\documentclass[FileMain.tex]{subfiles}
\gdef\sophong{Trường Mầm non, Tiểu học, THCS \& THPT}
\gdef\truong{ĐỀ ÔN TẬP HÈ}
\gdef\monhoc{Khoa học tự nhiên 6 (Ôn hè lớp 5 lên lớp 6)}
\gdef\ngaykt{.../.../2026}
\gdef\nh{2025 - 2026}
\gdef\thoigian{45}
\gdef\made{101}
\setcounter{section}{0}
%\tatloigiai      % <-- BỎ DẤU % Ở ĐẦU DÒNG NÀY để in BẢN HỌC SINH (ẩn lời giải)
%\hienthiloigiai
%\dongkeloigiai
\begin{document}
\section[Đề ôn tập hè Khoa học tự nhiên 6 - Mã đề \made]{Đề ôn tập hè Khoa học tự nhiên 6}

%%%==============Phần trắc nghiệm nhiều lựa chọn==============%%%
\subsection{Bài tập trắc nghiệm nhiều lựa chọn}\textit{\large Thí sinh trả lời từ câu 1 đến câu 12. Mỗi câu thí sinh chỉ chọn một phương án.}
\Opensolutionfile{ansex}[Ans/LGEX-DE_ON_HE_KHTN6_TONGHOP_MADE101]
\Opensolutionfile{ans}[Ans/Ans-DE_ON_HE_KHTN6_TONGHOP_MADE101]

%%%=============EX_1=============%%%
\begin{ex}%[6K1N1-1]
Khí nào sau đây chiếm khoảng $78\%$ thể tích không khí?
\choice
{Oxygen}
{Carbon dioxide}
{\True Nitrogen}
{Hơi nước}
\loigiai{
	Trong không khí, nitrogen chiếm khoảng $78\%$ thể tích, oxygen chiếm khoảng $21\%$, phần còn lại (khoảng $1\%$) là carbon dioxide, hơi nước và một số khí khác.
}
\end{ex}

%%%=============EX_2=============%%%
\begin{ex}%[6K1H1-2]
Phơi quần áo ướt ngoài nắng, sau vài giờ quần áo khô. Nước trong quần áo đã biến đổi như thế nào?
\choice
{Nước đông đặc}
{Nước ngưng tụ}
{\True Nước bay hơi vào không khí}
{Nước thấm ngược trở lại vào vải}
\loigiai{
	Dưới ánh nắng, nước trong quần áo (thể lỏng) nhận nhiệt và chuyển thành hơi nước (thể khí) bay vào không khí. Đó là sự bay hơi, nhờ vậy quần áo khô dần.
}
\end{ex}

%%%=============EX_3=============%%%
\begin{ex}%[6K1V1-3]
Ở vùng núi cao, nước thường sôi ở nhiệt độ thấp hơn $100^\circ$C. Nguyên nhân chủ yếu là
\choice
{nước ở vùng núi không tinh khiết}
{không khí trên núi lạnh hơn}
{\True áp suất khí quyển trên cao thấp hơn nên nhiệt độ sôi của nước giảm}
{nước trên núi chảy nhanh hơn}
\loigiai{
	Càng lên cao, áp suất khí quyển càng giảm. Áp suất trên mặt thoáng giảm làm cho nước sôi ở nhiệt độ thấp hơn $100^\circ$C. Vì vậy ở vùng núi cao nước sôi nhanh nhưng nhiệt độ sôi lại thấp hơn ở đồng bằng.
}
\end{ex}

%%%=============EX_4=============%%%
\begin{ex}%[6K2N1-1]
Trường hợp nào sau đây là một dung dịch?
\choice
{Cát được khuấy đều trong nước}
{Dầu ăn được khuấy với nước}
{\True Muối ăn tan hoàn toàn trong nước}
{Gạo trộn lẫn với đỗ xanh}
\loigiai{
	Dung dịch là hỗn hợp đồng nhất của chất tan và dung môi. Muối ăn tan hoàn toàn và phân bố đều trong nước tạo thành nước muối -- một dung dịch. Cát, dầu ăn không tan trong nước; gạo và đỗ chỉ trộn lẫn nên chỉ là hỗn hợp.
}
\end{ex}

%%%=============EX_5=============%%%
\begin{ex}%[6K2H1-2]
Để tách muối ăn ra khỏi dung dịch nước muối, ta nên làm cách nào sau đây?
\choice
{Dùng nam châm để hút muối}
{Dùng giấy lọc để lọc muối}
{\True Đun nóng dung dịch cho nước bay hơi, muối kết tinh ở lại}
{Để yên cho muối tự lắng xuống đáy}
\loigiai{
	Muối tan trong nước nên không lọc hay lắng được. Khi đun nóng, nước bay hơi còn muối không bay hơi sẽ kết tinh và ở lại, nhờ đó tách được muối ra khỏi dung dịch.
}
\end{ex}

%%%=============EX_6=============%%%
\begin{ex}%[6K2V1-3]
Hỗn hợp nào sau đây có thể tách các thành phần bằng nam châm?
\choice
{Muối và nước}
{Đường và nước}
{\True Mạt sắt và cát}
{Dầu ăn và nước}
\loigiai{
	Sắt bị nam châm hút còn cát thì không. Vì vậy có thể dùng nam châm để hút mạt sắt ra, tách nó khỏi cát. Các hỗn hợp còn lại không có thành phần bị nam châm hút.
}
\end{ex}

%%%=============EX_7=============%%%
\begin{ex}%[6K3N1-1]
Chất ở trạng thái khí có đặc điểm nào sau đây?
\choice
{Có hình dạng và thể tích xác định}
{Có hình dạng xác định nhưng thể tích thay đổi}
{\True Không có hình dạng xác định, chiếm toàn bộ không gian của vật chứa}
{Chỉ chiếm phần đáy của vật chứa}
\loigiai{
	Chất khí không có hình dạng và thể tích xác định; nó luôn lan tỏa và chiếm đầy toàn bộ không gian bên trong vật chứa.
}
\end{ex}

%%%=============EX_8=============%%%
\begin{ex}%[6K3H1-2]
Khi đun nóng, một mẩu nến rắn chảy ra thành nến lỏng. Đây là sự chuyển trạng thái
\choice
{lỏng sang rắn}
{\True rắn sang lỏng}
{rắn sang khí}
{lỏng sang khí}
\loigiai{
	Khi được đun nóng tới nhiệt độ thích hợp, nến từ thể rắn chuyển sang thể lỏng (nóng chảy). Đây là sự biến đổi trạng thái rắn sang lỏng, không tạo thành chất mới.
}
\end{ex}

%%%=============EX_9=============%%%
\begin{ex}%[6K4N1-1]
Sự biến đổi hóa học là sự biến đổi
\choice
{chỉ làm thay đổi hình dạng của chất}
{chỉ làm thay đổi trạng thái của chất}
{\True có sự tạo thành chất mới}
{làm cho chất bị chia nhỏ ra}
\loigiai{
	Dấu hiệu cơ bản của biến đổi hóa học là có chất mới được tạo thành, thường nhận biết qua sự thay đổi màu sắc, mùi vị hoặc tính tan của chất.
}
\end{ex}

%%%=============EX_10=============%%%
\begin{ex}%[6K4H1-2]
Hiện tượng nào sau đây là một biến đổi hóa học?
\choice
{Nước đá tan thành nước}
{Bẻ đôi một viên phấn}
{\True Đinh sắt để lâu trong không khí ẩm bị gỉ}
{Hòa tan đường vào nước}
\loigiai{
	Đinh sắt bị gỉ là do sắt tác dụng với các chất trong không khí ẩm tạo thành chất mới (gỉ sắt) có màu nâu đỏ, giòn, dễ gãy -- đó là biến đổi hóa học. Ba hiện tượng còn lại chỉ thay đổi trạng thái hoặc hình dạng, không tạo chất mới.
}
\end{ex}

%%%=============EX_11=============%%%
\begin{ex}%[6K5N1-1]
Dụng cụ nào sau đây dùng để đo chính xác thể tích chất lỏng?
\choice
{Cốc thủy tinh}
{\True Ống đong}
{Đũa thủy tinh}
{Kẹp ống nghiệm}
\loigiai{
	Ống đong có vạch chia thể tích nên dùng để đo chính xác thể tích chất lỏng. Cốc thủy tinh chỉ để đựng, đũa thủy tinh để khuấy, kẹp ống nghiệm để giữ ống khi đun.
}
\end{ex}

%%%=============EX_12=============%%%
\begin{ex}%[6K5H1-2]
Việc làm nào sau đây KHÔNG đúng quy tắc an toàn trong phòng thực hành?
\choice
{Đeo kính bảo hộ khi làm thí nghiệm}
{\True Nếm thử hóa chất để biết đó là chất gì}
{Đọc kỹ nhãn hóa chất trước khi sử dụng}
{Rửa tay sạch bằng xà phòng sau khi thực hành}
\loigiai{
	Tuyệt đối không được nếm hoặc ngửi trực tiếp hóa chất vì có thể gây ngộ độc, bỏng. Các việc còn lại đều là quy tắc an toàn cần thực hiện.
}
\end{ex}

\Closesolutionfile{ans}
\Closesolutionfile{ansex}
\bangdapan{Ans-DE_ON_HE_KHTN6_TONGHOP_MADE101}

%%%==============Phần trắc nghiệm đúng sai==============%%%
\subsection{Trắc nghiệm đúng sai}\textit{\large Thí sinh trả lời từ câu 1 đến câu 2. Trong mỗi ý a), b), c), d) ở mỗi câu thí sinh chọn đúng hoặc sai.}
\Opensolutionfile{ansex}[Ans/LGTF-DE_ON_HE_KHTN6_TONGHOP_MADE101]
\Opensolutionfile{ansbook}[Ansbook/AnsTF-DE_ON_HE_KHTN6_TONGHOP_MADE101]
\Opensolutionfile{ans}[Ans/Tempt-DE_ON_HE_KHTN6_TONGHOP_MADE101]
\setcounter{ex}{0}

%%%=============TF_1=============%%%
\begin{ex}%[6K1H1-4]
Cho các phát biểu sau về nước, không khí và đất:
\choiceTF
{\True Nước có thể tồn tại ở cả ba thể: rắn, lỏng và khí}
{Không khí có màu xanh nhạt và có vị ngọt}
{\True Mùn trong đất cung cấp chất dinh dưỡng cho cây trồng}
{Nitrogen chiếm khoảng $21\%$ thể tích không khí}
\loigiai{
	\begin{itemchoice}[T1,F2,T3,F4]
		\itemch Nước tồn tại ở ba thể: rắn (nước đá), lỏng (nước) và khí (hơi nước). Phát biểu đúng.
		\itemch Không khí trong suốt, không màu, không mùi, không vị. Phát biểu sai.
		\itemch Mùn được hình thành từ xác động, thực vật phân hủy, chứa nhiều chất hữu cơ và dinh dưỡng cung cấp cho cây. Phát biểu đúng.
		\itemch Nitrogen chiếm khoảng $78\%$, còn oxygen mới chiếm khoảng $21\%$ thể tích không khí. Phát biểu sai.
	\end{itemchoice}
}
\end{ex}

%%%=============TF_2=============%%%
\begin{ex}%[6K4H1-3]
Cho các phát biểu sau về sự biến đổi của chất:
\choiceTF
{\True Đốt cháy một tờ giấy thành tro là biến đổi hóa học}
{Nước đá tan thành nước lỏng là biến đổi hóa học}
{\True Có thể nhận biết biến đổi hóa học qua dấu hiệu đổi màu, đổi mùi hoặc đổi tính tan}
{\True Thức ăn để lâu bị ôi thiu là một biến đổi hóa học}
\loigiai{
	\begin{itemchoice}[T1,F2,T3,T4]
		\itemch Đốt giấy tạo ra tro -- chất mới khác hẳn giấy ban đầu, nên là biến đổi hóa học. Phát biểu đúng.
		\itemch Nước đá tan thành nước chỉ thay đổi trạng thái (rắn sang lỏng), không tạo chất mới, là biến đổi vật lí. Phát biểu sai.
		\itemch Biến đổi hóa học thường được nhận biết qua sự thay đổi màu sắc, mùi vị, tính tan. Phát biểu đúng.
		\itemch Thức ăn ôi thiu đã sinh ra chất mới có mùi, vị khác, nên là biến đổi hóa học. Phát biểu đúng.
	\end{itemchoice}
}
\end{ex}

\Closesolutionfile{ans}
\Closesolutionfile{ansbook}
\Closesolutionfile{ansex}
\bangdapanTF{AnsTF-DE_ON_HE_KHTN6_TONGHOP_MADE101}

%%%==============Phần bài tập trả lời ngắn==============%%%
\subsection{Bài tập trả lời ngắn}\textit{\large Thí sinh trả lời từ câu 1 đến câu 3.}
\Opensolutionfile{ansex}[Ans/LGSA-DE_ON_HE_KHTN6_TONGHOP_MADE101]
\Opensolutionfile{ansexh}[Ans/AnsSA-DE_ON_HE_KHTN6_TONGHOP_MADE101]
\setcounter{ex}{0}

%%%=============SA_1=============%%%
\begin{ex}%[6K1N1-5]
Ở điều kiện thường, nước sôi ở nhiệt độ bao nhiêu độ C?
\shortans{$100$}
\loigiai{
	Ở điều kiện thường (áp suất khí quyển chuẩn), nước sôi ở $100^\circ$C.
}
\end{ex}

%%%=============SA_2=============%%%
\begin{ex}%[6K1N1-6]
Nitrogen chiếm khoảng bao nhiêu phần trăm thể tích không khí?
\shortans{$78$}
\loigiai{
	Nitrogen chiếm khoảng $78\%$ thể tích không khí.
}
\end{ex}

%%%=============SA_3=============%%%
\begin{ex}%[6K1N1-7]
Ở điều kiện thường, nước đông đặc (hóa thành nước đá) ở nhiệt độ bao nhiêu độ C?
\shortans{$0$}
\loigiai{
	Ở điều kiện thường, nước đông đặc thành nước đá ở $0^\circ$C.
}
\end{ex}

\Closesolutionfile{ansexh}
\Closesolutionfile{ansex}
\bangdapanSA{AnsSA-DE_ON_HE_KHTN6_TONGHOP_MADE101}

%%%==============Phần bài tập tự luận==============%%%
\subsection{Bài tập tự luận}\textit{\large Thí sinh trả lời từ bài 1 đến bài 2.}
\Opensolutionfile{ansbth}[Ans/LGBT-DE_ON_HE_KHTN6_TONGHOP_MADE101]
\Opensolutionfile{ansbt}[Ans/AnsBT-DE_ON_HE_KHTN6_TONGHOP_MADE101]

%%%=============BT_1=============%%%
\begin{bt}%[6K4V1-1]
Phân biệt sự khác nhau giữa biến đổi vật lí và biến đổi hóa học. Mỗi loại cho hai ví dụ minh họa.
\loigiai{
	\begin{itemize}
		\item \textbf{Biến đổi vật lí:} chất chỉ thay đổi hình dạng, kích thước hoặc trạng thái mà \textbf{không} tạo thành chất mới.\\
		Ví dụ: xé nhỏ một tờ giấy; nước đá tan thành nước.
		\item \textbf{Biến đổi hóa học:} có sự \textbf{tạo thành chất mới}, thường nhận biết qua sự đổi màu, đổi mùi hoặc đổi tính tan.\\
		Ví dụ: đốt giấy thành tro; đinh sắt để lâu trong không khí ẩm bị gỉ.
	\end{itemize}
}
\end{bt}

%%%=============BT_2=============%%%
\begin{bt}%[6K2V1-4]
a) Trình bày cách tách muối ăn ra khỏi dung dịch nước muối và giải thích vì sao làm được như vậy.

b) Khi đun nóng hóa chất trong ống nghiệm, vì sao phải hướng miệng ống nghiệm về phía không có người?
\loigiai{
	\textbf{a)} Đun nóng dung dịch nước muối. Khi đun, nước bay hơi dần, còn muối không bay hơi nên kết tinh và ở lại, nhờ đó thu được muối ăn.\\
	Làm được như vậy vì muối tan trong nước tạo thành dung dịch; khi nước bay hơi hết thì muối tách ra ở dạng tinh thể.

	\textbf{b)} Vì khi đun, chất lỏng trong ống nghiệm có thể sôi mạnh và bắn ra ngoài qua miệng ống. Hướng miệng ống về phía không có người để tránh hóa chất nóng bắn vào mình hoặc người xung quanh, bảo đảm an toàn.
}
\end{bt}

\Closesolutionfile{ansbt}
\Closesolutionfile{ansbth}

\begin{center}
 \rule[4pt]{2cm}{1pt}\,\large\bfseries Hết\,\rule[4pt]{2cm}{1pt}
\end{center}
\label{x}
\end{document}
